\clearpage
\begingroup
\fancyhead[L]{\small THEORIE DER BEWEGUNG DES MONDES.}
\renewcommand{\headrulewidth}{0pt}
\SourceLeaf{626}
\section*{DRITTER ABSCHNITT.}
\begin{center}
\itshape Erste Annäherung.
\end{center}

\par\noindent\hspace*{\parindent}Zuerst ist hier \(m''\) zu bestimmen.
Aus der Theorie der elliptischen Bewegung ist bekannt, dass die folgende Gleichung gilt, wenn man die
Bewegung der Erde um die Sonne allein als Ergebnis der Anziehungskräfte dieser
beiden Himmelskörper betrachtet und die mittlere Bewegung der Sonne \(=nu\) sowie
[die] mittlere Entfernung der Erde von der Sonne \(=B\) setzt.
\[
m'+m''=nnB^3\frac{\mathrm{d}u^2}{\mathrm{d}t^2};
\]
\noindent Setzt man also \(\dfrac{m'}{m''}=\mu\), wobei dieses Verhältnis hinreichend genau
bekannt ist \(\left(=\dfrac{1}{365412}\right)\), so erhält man
\[
m''=nnB^3\frac{\mathrm{d}u^2}{\mathrm{d}t^2}\frac{1}{1+\mu}
\mathpunct{\raisebox{.4ex}{.}}
\]
\noindent Ferner ist
\[
R'=\BareRoot{\left(RR-\frac{2R}{p}\cos(v-V)+\frac{1+\theta\theta}{pp}\right)}
\]
\noindent oder, wenn man die höheren Potenzen von \(\dfrac{1}{R}\) und \(\theta\) weglässt,
\[
\begin{aligned}
\frac{1}{R'^3}&=\frac{1}{R^3}\left\{1+\frac{3\cos(v-V)}{Rp}\right.\\
&\qquad\left.+\frac{15}{4}\frac{1}{RRpp}(1+\cos2(v-V))
-\frac32\frac{1}{RRpp}\right\}\mathpunct{\raisebox{.4ex}{.}}
\end{aligned}
\]
\noindent Hieraus folgt
\[
\begin{aligned}
\pi&=\frac{nn}{1+\mu}\frac{B^3}{R^3}\left\{\frac32\sin2(v-V)\right.\\
&\qquad\left.+\frac{1}{Rp}\left(\frac38\sin(v-V)+\frac{15}{8}\sin3(v-V)\right)\right\}\\[5pt]
\omega&=\frac{nn}{1+\mu}\frac{B^3}{R^3}\left\{\frac32+\frac32\cos2(v-V)\right.\\
&\qquad\left.+\frac{1}{Rp}\left(\frac{33}{8}\cos(v-V)+\frac{15}{8}\cos3(v-V)\right)\right\}\\[5pt]
\psi&=\frac{nn}{1+\mu}\frac{B^3}{R^3}\left\{\frac12+\frac32\cos2(v-V)\right.\\
&\qquad\left.+\frac{1}{Rp}\left(\frac98\cos(v-V)+\frac{15}{8}\cos3(v-V)\right)\right\}
\mathpunct{\raisebox{.4ex}{.}}
\end{aligned}
\]
\SourceLeaf{627}
\noindent Setzt man schließlich die Sonnenexzentrizität \(=\varepsilon\) und die mittlere
Anomalie \(=a\), so erhält man
\[
\frac{R}{B}=1+\tfrac12\varepsilon\varepsilon+\varepsilon\cos a
-\tfrac12\varepsilon\varepsilon\cos2a,
\]
\noindent also
\[
\frac{B^3}{R^3}=1+\tfrac32\varepsilon\varepsilon-3\varepsilon\cos a
+\tfrac92\varepsilon\varepsilon\cos2a,
\]
\noindent und
\[
V=nu-2\varepsilon\sin a+\tfrac54\varepsilon\varepsilon\sin2a.
\]

\par\noindent\hspace*{\parindent}Bei der Einsetzung der elliptischen Bewegung
in den Termen \(\sin\lambda\) und \(\cos M\) müssen die Zeichen \(\lambda\), \(M\)
unbestimmt bleiben. Die Einsetzung von \(v-\mathit{Const.}\)
für diese Zeichen ist unzulässig. Die Störung durch die Sonne könnte
beim Knoten und beim Apogäum eine Bewegung mit einem nichtperiodischen Anteil hervorrufen,
sodass \(\lambda\) und \(M\) die Form \(\zeta v-\mathit{Const.}\) annähmen,
wobei sich \(\zeta\) nur wenig von der Einheit unterscheiden würde.
In diesem Fall wäre die Annahme der elliptischen Bewegung völlig falsch,
wenn man für \(\lambda\) von \(=\zeta v-\mathit{Const.}\) zu \(=v-\mathit{Const.}\) überginge,
weil sich diese Abweichung immer weiter summieren würde.
Wir lassen die Argumente \(\lambda\) und \(M\) daher unbestimmt. Genauer gesagt
setzen wir \(\lambda=\zeta v-\mathit{Const.}\), \(M=\eta v-\mathit{Const.}\),
wobei \(\zeta\), \(\eta\) Koeffizienten sind, deren Wert ohne die Störung
durch die Sonne \(=1\) wäre. Er bleibt unbestimmt,
bis die Rechnung ihn ergibt.

\par\noindent\hspace*{\parindent}Lassen wir bei dieser ersten Annäherung in der dritten Gleichung die
Terme weg, die \(\dfrac{nn}{R}\), \(nn\varepsilon\) und \(nne\) enthalten,
so ergibt sich
\[
\begin{aligned}
0&=\frac{\mathrm{dd}\theta}{\mathrm{d}v^2}+(1+\tfrac32nn)\theta\\
&\qquad+\tfrac32nn\left(\theta\cos2(v-V)
-\frac{\mathrm{d}\theta}{\mathrm{d}v}\sin2(v-V)\right);
\end{aligned}
\]
\noindent Für \(V\) dürfen wir hier \(nv[+\mathit{Const.}]\) einsetzen.
Schreiben wir folglich für \(v-nv[-\mathit{Const.}]\) das Zeichen \(E\) und für
\(\theta\) den Ausdruck \(g\sin\lambda\), so ergibt sich
\[
0=\frac{\mathrm{dd}\theta}{\mathrm{d}v^2}+(1+\tfrac32nn)\theta
-\tfrac32nng\sin(2E-\lambda).
\]

\par\noindent\hspace*{\parindent}Diese Gleichung zeigt unmittelbar, dass der Ansatz
\(\lambda=v-\mathit{Const.}\) durch
\(=v\BareRoot{(1+\tfrac32nn)}-\mathit{Const.}\) zu ersetzen ist.
Also ist \(\dfrac{\mathrm{d}\lambda}{\mathrm{d}v}=1+\tfrac34nn\). Der
Knoten hat damit eine rückläufige Bewegung, die sich zur mittleren Bewegung
des Mondes wie \(\tfrac34nn:1\) verhält. Daraus folgt der erste angenäherte
\SourceLeaf{628}Wert
\[
\theta=g\sin\lambda+(\tfrac38n+\tfrac{3}{32}nn)g\sin(2E-\lambda).
\]

\par\noindent\hspace*{\parindent}Um zu sehen, welche numerischen Koeffizienten diese erste
Annäherung ergibt, verwenden wir die folgenden Beobachtungswerte.
\[
\begin{array}{r@{\;}c@{\;}r}
\log n&=&8{,}873\,9096\\
\log nn&=&7{,}747\,8192\\
\log\tfrac34&=&9{,}875\,0613\\\cline{2-3}
&&7{,}622\,8805
\end{array}
\]
\noindent Somit ist \(\tfrac34nn=0{,}004\,1964\,35\), während die Beobachtungen ergeben
\(\dfrac{\mathrm{d}\lambda}{\mathrm{d}v}=1{,}004\,0218\,706\).

\par\noindent\hspace*{\parindent}Ferner ist \(g\) in Bogensekunden \(=18567\overset{\prime\prime}{,}933\).
\[
\begin{array}{r@{\;}c@{\;}r@{\qquad}r@{\;}c@{\;}r}
\log&=&4{,}268\,7636&&&4{,}268\,7636\\
\log\tfrac38n&=&8{,}447\,9409&\log\tfrac{3}{32}nn&=&6{,}719\,7905\\\cline{2-3}\cline{5-6}
&&2{,}716\,7045&&&0{,}988\,5541\\[3pt]
&&520\overset{\prime\prime}{,}840&&&+9\overset{\prime\prime}{,}740
\end{array}
\]
\noindent Der Koeffizient ist also \(=530\overset{\prime\prime}{,}580\) und stimmt nach einer
später zu erklärenden Reduktion genau mit den Beobachtungen überein.
\clearpage
\endgroup
